\subsection{From the plane branch to the actual train output}
\label{sec:canonical-transfer}

All cell measures in this subsection are unnormalized. Put
\[
\begin{gathered}
 \Theta=\nu^{-1},\quad c=\nu^{-2},\quad e=(\nu c^2)^{-1}=\nu^3,\quad
 a=1+\nu s,\\
 \theta=(\nu a)^{-1},\quad d=c/\sqrt a=\sqrt{\theta/e},\quad
 N=e^{\lambda/\nu}\nu^{-\kappa}.
\end{gathered}
\]
Write $P_d$ for periodization, $K_d$ for cylindrical Biot--Savart,
and $X_{m,d}$ for the centrally even zero-circulation class with
weight $\langle y\rangle^{2m}$ on the centered cell, $m>8$.
Its first moments vanish. The canonical plane branch is that of
Proposition~\ref{prop:short-canonical}; set
$M(\theta)=\mathfrak M(V_{\rm can}(\theta))$. Angular integration
and cell summation give $Q_dP_dV_{\rm can}(\theta)=M(\theta)$.

\paragraph{The actual residual, including moving images.}
For $z=\Theta-\theta$, the identity $d_z/d=-1/(2\theta)$ gives
\begin{equation}\label{eq:ct-moving-period}
 \theta^{-1}D(P_df)-\partial_z(P_df)
       =\theta^{-1}P_dDf
\end{equation}
when $f$ is independent of $z$. Thus periodization commutes with
the combined diffusion/dilation/moving-cell derivative, not with
each term separately. Reynolds' boundary term cancels the dilation
seam term. All calculations below may first be made on the fixed
physical cylinder and then rescaled.

For completeness the improved even image error is elementary.
The symmetrically summed planar kernels satisfy
\[
 K_d(\xi)-K_\infty(\xi)
      =d^{-2}B_1\xi+O(d^{-4}|\xi|^3),\qquad |\xi|\le3d/8.
\]
Indeed the opposite translates cancel the constant and quadratic
Taylor coefficients, while their $j$th derivatives are
$O((|n|d)^{-1-j})$. The affine convolution vanishes because both
mass and first moments are zero. On $|y|\le d/4$, $|x|\le d/8$,
the remaining convolution is bounded by
$Cd^{-4}(1+|x|^3)\|f\|_{X_m}$. The excluded source tail is at most
$Cd^{1-m}\|f\|_{X_m}$ in $L^1$ and its far kernel costs $Cd^{-1}$.
Multiplication by the Gaussian observation coefficient controls
the polynomial in $x$; its complement $|x|>d/8$ is exponentially
small. The same reasoning after $d\partial_d$ gives the same bound.
Consequently the localized linearized interaction error is
$O(d^{-4})$ on this even class. For a periodized Gaussian input,
cross transport between different cells is exponentially small,
using its fixed Gaussian derivative norms.

Let $E_c$ be the actual train vorticity and
$h=a^2E_c(s,\sqrt a\,y)-P_dE_0$.
The train comparison \eqref{sct:comparison}, after rescaling, gives
\[
 \|\partial^j h\|_2+
 \|\partial^jK_dh\|_{2\cap\infty}
       \le C_j\theta e^{-b_jd^2}.
\]
Here the curl estimate uses one additional derivative. Since
$\theta=ed^2$, the right side is at most $C_je e^{-b'_jd^2}$.
There is no reliance on smallness of an unbounded differential
operator. To add fixed polynomial weights, the heat-coordinate
vorticity energy with weight $\langle y\rangle^{4m}$ and finitely
many derivatives gives the crude bound $P_\nu e^{C_{m,j}/\nu}$:
transport coefficients and their selected derivatives are uniformly
bounded by the preceding unweighted comparison, their clock interval
has length at most $\nu^{-1}$, and the diffusion/dilation contribution
integrates $C_{m,j}\theta^{-1}$, costing only a power of $\nu$.
The periodic seam diffusion term is nonpositive and the moving seam
cancels as in \eqref{eq:ct-moving-period}. Interpolating that bound
with the unweighted error, and using $d^2\ge r/e$ on any fixed
$[r,\Theta]$, absorbs $e^{C/\nu}$ and every fixed polynomial into
a smaller Gaussian for sufficiently small $\nu$:
\begin{equation}\label{eq:ct-base-error}
 \|h\|_{H_m^j}+\|K_dh\|_{W^{j,\infty}}
       \le C_{m,j}e\,e^{-b_{m,j}\theta/e},
 \qquad r\le\theta\le\Theta.
\end{equation}

Let $w=\mathcal U_{\nu,c}(s,0)P_cq$ be the actual bare column and
write $w(s,x)=Na^{-2}Z(\theta,x/\sqrt a)$.
The canonical asymptotic gives
\[
 \|e^{\lambda\Theta}\Theta^\kappa
       [Z(\Theta)-P_cV_{\rm can}(\Theta)]\|_{X_{m,c}}\le C\nu.
\]
Applying
\eqref{eq:ct-moving-period}, the image estimate, and
\eqref{eq:ct-base-error} to $V_d=P_dV_{\rm can}$ gives the following
residual in the equation $-\partial_\theta Z=\mathcal L_\theta Z$:
\begin{align}
 \|e^{\lambda\theta}\theta^\kappa
       (-\partial_\theta-\mathcal L_\theta)V_d\|_{X_{m,d}}
 &\le C_b(e^2\theta^{-2}+e\,e^{-a_0\theta/e}),
 &&b\le\theta\le\Theta, \label{eq:ct-long-residual}\\
 \|(-\partial_\theta-\mathcal L_\theta)V_d\|_{X_{m,d}}
 &\le C_b(e^2\theta^{-2}+e\,e^{-a_0\theta/e}),
 &&r\le\theta\le b. \label{eq:ct-late-residual}
\end{align}
The image term needs only the Gaussian zeroth norm of the branch;
cross transport and the base-error terms use its fixed derivative
budgets. Their polynomial factors in $1+\theta$ are absorbed into
the displayed Gaussian. This explicitly accounts for the actual
nonlinear train, rather than substituting its independent-vortex sum.

The full-input forward bound \eqref{eq:short-causal-long-forward}
and Duhamel, including this normalized inlet error, imply
\begin{equation}\label{eq:ct-long-comparison}
 \|e^{\lambda\theta}\theta^\kappa(Z-V_d)(\theta)\|_{X_{m,d}}
 \le C_b\{\nu+e^2/\theta+e^{-a_1\theta/e}\},\qquad b\le\theta\le\Theta.
\end{equation}
In particular the error at $b$ is $O_b(\nu)$, and both columns
have an exponentially decaying, polynomially weighted majorant
for $\theta\ge b$. No terminal projection or backward parabolic
propagator has been used.

\paragraph{Kinetic and tested moment continuation through overlap.}
The embeddings
\begin{equation}\label{eq:ct-embeddings}
 |Q_d f|+\|K_df\|_2\le C_m\|f\|_{X_{m,d}},\qquad
 Q_df=-\frac2\pi\int\eta^2f,\qquad d\ge1
\end{equation}
are uniform in $d$. The moment bound is weighted Cauchy--Schwarz.
For velocity, zero mass gives
$|\widehat f(\xi)|\le|\xi|\int|y||f(y)|dy$ at low cylindrical
frequencies. In Plancherel, the $O(d)$ horizontal frequencies
of modulus at most one cancel the measure factor $d^{-1}$;
the high frequencies are bounded by $\|f\|_2$. This proves
\eqref{eq:ct-embeddings} without a factor $1+d$.

Let $\mathcal U_Z(\theta,r)$ be the exact heat-coordinate evolution.
The velocity energy inequality and \eqref{sct:integrated} give,
for $0<\theta\le r\le b$,
\begin{equation}\label{eq:ct-kinetic-impulse}
 \|K_{d(\theta)}\mathcal U_Z(\theta,r)h\|_2
       \le C_b(r/\theta)\|K_{d(r)}h\|_2.
\end{equation}
Indeed the input physical kinetic norm is
$\nu r\|K_{d(r)}h\|_2$ and the integrated strain costs at most
$C(r+e)$. The exact physical quadrupole law is
\[
 \frac{d}{ds}Q_c(w)
    =-\frac4\pi\int(v_1u_2+v_2u_1)
    =-\frac4\pi\int(\widetilde v_1\widetilde u_2+
                              \widetilde v_2\widetilde u_1).
\]
The mean terms vanish because the mean vertical velocities vanish.
Since $Q$ is invariant under heat scaling, the same identity gives
a tested bound uniform in $\nu$. Before overlap
$\|\widetilde v(s(\theta))\|_2\le C\nu\theta$;
after overlap it is at most
$C\nu e\exp[-a(e/\theta-1)]$. Thus at every age
$\nu^{-1}\theta^{-2}\|\widetilde v(s(\theta))\|_2\le C/\theta$,
and
\begin{align}
 |Q\mathcal U_Z(\theta,r)h|
 &\le |Qh|+C_b r\log(r/\theta)\|K_{d(r)}h\|_2,
                      \label{eq:ct-partial-moment}\\
 \int_0^r\theta^{p-1}|Q\mathcal U_Z(\theta,r)h|\,d\theta
 &\le C_{b,p}r^p(|Qh|+r\|K_{d(r)}h\|_2),\quad p>0.
                      \label{eq:ct-tested-impulse}
\end{align}
The second follows from $\int_0^r\theta^{p-1}\log(r/\theta)d\theta=r^p/p^2$.
Similarly the absolute stress, divided by its scale
$\nu^{3+2\Re\mu}/c$, for an impulse at $r$ is bounded by
\begin{equation}\label{eq:ct-absolute-impulse}
 C_{K,b}r^{p+1}\max\{1,(e/r)^{1/4}\}\|K_{d(r)}h\|_2,
 \qquad p=\Re\mu+3/2.
\end{equation}
For $\theta\ge e$ insert $\|v\|_2\le C\nu\theta$ and
$\|u\|_2\le C_b\nu r\|Kh\|_2$. For $\theta\le e$ insert
$\|v\|_2\le C\nu e^{1/4}\theta^{3/4}$ instead. Its integral is
$Cr e^{1/4}\min(e,r)^{p-1/4}\|Kh\|_2$, proving the assertion.
All these identities first hold for smooth data and sources.
Duhamel and the integrable displayed bounds give the stated
completion on the fixed physical cylinder; the actual decaying
mean stream is retained when identifying its quadrupole.

For each fixed $r>0$, apply \eqref{eq:ct-kinetic-impulse} and
\eqref{eq:ct-partial-moment} to the residual and the error at $b$.
The kinetic residual integral is bounded by
$C_be^2\theta^{-1}\log(b/\theta)$; the partial-moment factors are
bounded on $r\le\theta\le b$. Thus
\begin{equation}\label{eq:ct-compact-comparison}
 \sup_{r\le\theta\le b}
 \bigl(\|K_d(Z-V_d)(\theta)\|_2+
                |Q_d(Z-V_d)(\theta)|\bigr)\longrightarrow0.
\end{equation}
The plane bounds and periodization imply
$\limsup_{\nu\to0}(|Q_dZ(r)|+\|K_dZ(r)\|_2)\le C_b$ uniformly
in $0<r\le b$. Apply the impulse estimates to the homogeneous actual
column with data $Z(r)$, not to its discrepancy from $V_d$.
For $p$ in a compact subset of $(1/4,\infty)$,
\[
 \limsup_{\nu\to0}\int_0^r\theta^{p-1}|Q_dZ(\theta)|d\theta
       \le C_br^p,
\]
and \eqref{eq:ct-absolute-impulse}, with $e<r$, bounds its normalized
absolute stress tail by $C_br^{p+1}$. Letting first $\nu\to0$ and
then $r\downarrow0$ therefore gives the required tested convergence
through the entire overlap tail.

\paragraph{Mellin output and stress defects.}
Fix $K\Subset\{0<\Re\mu<1/4\}$, put $\sigma=\Re\mu$,
$u_o=\nu^{-3/2}$, and $\theta_o=(\sqrt\nu+\nu)^{-1}$.
For the actual column define
$A_+(s)=-\pi s^{-\mu}Q_c(w(s))/(4c)$.
The exact substitution
$s=(1-\nu\theta)/(\nu^2\theta)$ yields
\begin{align}
 I_A:=\int_{u_o}^{\infty}A_+(s)s^{-5/2}ds
 &=-\frac{\pi N}{4c}\nu^{3+2\mu}
   \int_0^{\theta_o}\theta^{\mu+1/2}
       (1-\nu\theta)^{-\mu-5/2}QZ(\theta)d\theta \notag\\
 &=-\frac{\pi N}{4c}\nu^{3+2\mu}
       [\mathcal J(\mu)+o_K(1)].
                                      \label{eq:ct-mellin}
\end{align}
Use \eqref{eq:ct-long-comparison} for the exponentially integrable
large-clock part and the compact-clock convergence and tested tail
bound above for the rest. The error is locally uniform and holomorphic.
An extra $s^{-k/2}$ costs
$\nu^k\theta^{k/2}(1-\nu\theta)^{-k/2}$, $k=0,1$.
At $u_o$, every selected state or stress observation divided by $N$
is $P_\nu e^{-\gamma'/\sqrt\nu}$ for any $\gamma'<\gamma$.
The partial-moment bound gives $QZ(\theta)=O_\nu(1+|\log\theta|)$
as $\theta\downarrow0$, so the upper weighted moment endpoint is zero.
If the old cutoff
starts at $\sigma_{\rm in}/\nu<u_o$, its entire preceding part
has the same exponentially small coefficient after polynomial
weights; that source and cutoff transition have not been removed.

For $u=K_cw$, $u_\mu=s^{-\mu}u$, set
\[
\begin{gathered}
 E_\eta=\frac2c\int v_1u_{\mu,1},\quad
 E_q=\frac2c\int v_2u_{\mu,2},\quad
 C=\frac1c\int(v_1u_{\mu,2}+v_2u_{\mu,1}),\\
 D_c=E_\eta+E_q,\quad R_c=E_q-iC,\quad
 \mathcal A_\nu=|N|\nu^{3+2\sigma}/c.
\end{gathered}
\]
These are bilinear pairings. Define the stronger absolute component
quantity
\[
 S_{\rm abs}(s)=c^{-1}\int
       |v\otimes u_\mu+u_\mu\otimes v|\,dq\,d\eta.
\]
For $k=0,1$,
\begin{equation}\label{eq:ct-absolute-stress}
 \int_{u_o}^\infty s^{-3/2-k/2}
       S_{\rm abs}(s)ds\le C_{K,k}\mathcal A_\nu\nu^k.
\end{equation}
This controls each absolute tensor component before cell integration,
and hence also $|E_\eta|+|E_q|+|C|$, up to an absolute constant.
Indeed the integrand before its time weight is at most
$C|N|c^{-1}s^{-\sigma}\|v\|_2\|u/N\|_2$.
For $\theta\ge b$, $\|u/N\|_2\le Ca^{-1}H(\theta)$ with
$H$ exponentially decreasing times a polynomial. Changing variables
gives the integrable majorant
$C\mathcal A_\nu\nu^k\theta^{\sigma+3/2+k/2}H(\theta)$.
For $\theta\le b$, energy gives $\|u/N\|_2\le C_b\nu$.
Before overlap $\|v\|_2\le C/(\nu s)$; after overlap
$\|v\|_2\le Cc^{-2}(s/s_c)^{-3/4}$.
The two integrals are bounded by
$C\mathcal A_\nu\nu^k$ and
$C\mathcal A_\nu\nu^ke^{\sigma+3/2+k/2}$ respectively.

For a plane radial background in positive angular sector two,
$E_\eta=-E_q$ and $E_q=iC$ by direct angular integration.
On each fixed positive clock interval,
\eqref{eq:ct-compact-comparison} and the actual background comparison
therefore make the scaled $D_c,R_c$ tend to zero. Their absolute
small-clock tails are bounded by \eqref{eq:ct-absolute-impulse} with
data $Z(r)$, hence by $C_K\mathcal A_\nu r^{\sigma+5/2}$ in the
$\nu\to0$ limit. Large-clock domination, followed by $r\downarrow0$,
proves
\begin{equation}\label{eq:ct-stress-defect}
 \int_{u_o}^{\infty}(D_c-3R_c)s^{-3/2}ds=o_K(\mathcal A_\nu).
\end{equation}
The unforced moment law is exactly $C=A_+'+\mu A_+/s$; hence
\begin{equation}\label{eq:ct-anisotropy}
 \int_{u_o}^{\infty}(2E_q-E_\eta)s^{-3/2}ds
   =3i(\mu+3/2)I_A+o_K(\mathcal A_\nu),
\end{equation}
including the boundary term
$3i[A_+s^{-3/2}]_{u_o}^{\infty}$, just shown negligible.
For the negative channel the sign of $i$ reverses and
$N\mathcal J(\mu)$ becomes
$\bar N\,\overline{\mathcal J(\bar\mu)}$.
